% !TEX root = ../main.tex
\section{Data Analysis}
\label{sec:data}

We show how our approach to data fairness can be used in practice by applying our estimators to real data.
We analyze the Adult-Income dataset and the Law School dataset, two datasets common in the fairness literature, and discuss the interpretations and implications of our results from a data fairness perspective.
To perform our analysis, we split the data into a 60/40 train/test split. For parity and equal opportunity, we train SuperLearner ensembles \cite{laanSuperLearner2007} for the outcome, group, and joint group--outcome distributions. For CMI, we fit a calibrated histogram gradient boosting classifier to the joint group--outcome distribution and derive its marginals from that fit.
For both datasets, our analysis seeks to determine to what extent the distribution of $Y$ given $X$ violates definitions of fairness.
We also report model-fairness intervals conditional on the fitted outcome model. Both targets use the same outcome predictions and held-out observations. The probabilistic data-fairness estimators include the nuisance corrections described above, whereas the model-fairness estimators compare group means of the fitted scores. With the same empirical influence-function variance convention, the thresholded estimates and intervals coincide numerically for the two targets.
Thresholded inference assumes the stability conditions in Appendix \ref{app:threshold_stability}, which we have not verified for these fits.

\begin{table}
  \centering
  \input{tables/real_data_inference}
  \caption{Data- and model-fairness estimates with nominal $95\%$ confidence intervals in parentheses. Parity and equal-opportunity comparisons share the fitted outcome model and evaluation split. CMI is reported only for data fairness; its intervals near zero should be interpreted with caution (Section \ref{sec:cmi}).}
  \label{tab:data_results}
\end{table}

\paragraph{Adult-Income.}
The Adult-Income dataset contains economic and demographic information on individuals living in the United States \cite{dingRetiringAdultNew2021} \cite{floodIntegratedPublicUse2020}.
For our analysis, the outcome is $Y = \ind\{\text{income} > 50,000\}$ and $G = \ind\{\text{is male}\}$.
We find that both parity and probabilistic parity are violated, with confidence intervals bounded away from zero.
Finally, the CMI estimate is near zero, but Wald inference is non-regular at conditional independence, so this result does not establish independence.

\paragraph{Law School.}
The Law School dataset contains educational, demographic, and economic information about law students \cite{wightmanLSACNationalLongitudinal1998} \cite{lequySurveyDatasetsFairnessaware2022}.
Here, our outcome of interest is $Y=\ind\{\text{passed the bar}\}$ and $G = \ind\{\text{is white}\}$.
Again we obtain positive estimates with intervals bounded away from zero for parity and equal opportunity.
We note that both equal opportunity estimates have narrower intervals than in the Adult dataset.
Positive outcomes are less prevalent in Adult than in Law, so conditioning on $Y=1$ retains a smaller fraction of the Adult observations.
This is consistent with the wider equal-opportunity intervals observed for Adult.
